\documentclass[10pt]{article}

\usepackage[utf8]{inputenc}
\usepackage[margin=1.0in]{geometry}
\usepackage{amsmath,amssymb,amsfonts}
\usepackage{booktabs}
\usepackage{graphicx}
\usepackage{url}
\usepackage{microtype}
\usepackage{hyperref}

\hypersetup{
    colorlinks=true,
    linkcolor=blue,
    citecolor=blue,
    urlcolor=blue
}

\title{Auditable Failure Memory for ROS 2 Navigation:\\Supplementary Materials \& Verification Audit Details}

\author{
    Author One\thanks{Department of Robotics, University Placeholder.}
    \and
    Author Two\thanks{Department of Computer Science, Institute Placeholder.}
    \and
    Author Three\thanks{Autonomous Systems Laboratory, University Placeholder.}
}
\date{}

\begin{document}

\maketitle

\section*{Document Overview}
This document provides supplementary technical specifications, precondition integrity constraints, cryptographic audit procedures, source-code traceability mappings, and kinematic ledger decomposition tables accompanying the paper \emph{Auditable Failure Memory for ROS 2 Navigation: An Exploratory Comparison with Spatial Caching}.

\section{S1. Precondition Integrity \& Doorway Clearance Details}
The passage state evaluation in \texttt{src/doorway\_evaluator.py} enforces strict sensor-TF alignment rules prior to emitting a clearance classification:

\subsection{S1.1 Precondition Rules}
\begin{enumerate}
    \item \textbf{Finite Transforms}: TF translation vector $\mathbf{t} = [x, y, z]^T$ and yaw angle $\psi$ must be finite real numbers ($-\infty < x, y, z, \psi < \infty$).
    \item \textbf{Temporal Alignment}: Scan and TF timestamps must satisfy $|t_{\text{tf}} - t_{\text{scan}}| \le 0.80\,\text{s}$.
    \item \textbf{Data Freshness}: Current simulation clock $t_{\text{sim}}$ must satisfy $t_{\text{sim}} - t_{\text{tf}} \le 1.0\,\text{s}$ and $t_{\text{sim}} - t_{\text{scan}} \le 1.5\,\text{s}$.
    \item \textbf{Ray Projections}:
    \begin{itemize}
        \item Number of finite laser scan ranges $N_{\text{valid}} > 0$.
        \item Number of projected rays intersecting the doorway bounding volume $N_{\text{intersect}} > 0$.
    \end{itemize}
    \item \textbf{Costmap Availability}: When costmap subgrids are provided, total cells in the Region of Interest (ROI) $N_{\text{total\_cells}} > 0$.
\end{enumerate}
If any precondition fails, the evaluator emits $\text{State}(B) = \text{UNKNOWN}$ and does not alter memory state.

\subsection{S1.2 Bounding Volume \& Threshold Parameters}
\begin{itemize}
    \item Doorway center: $x = 0.0\,\text{m}, y = 1.20\,\text{m}$.
    \item ROI bounding box: $x \in [-0.15, 0.15]\,\text{m}, y \in [0.95, 1.45]\,\text{m}$.
    \item Minimum pass-through rays for \texttt{FREE}: $N_{\text{pass\_through}} \ge 8$.
    \item Maximum hits inside for \texttt{FREE}: $N_{\text{hits\_inside}} = 0$.
    \item Hits threshold for \texttt{OCCUPIED}: $N_{\text{hits\_inside}} \ge 5$ or $N_{\text{costmap\_occupied\_cells}} > 0$.
\end{itemize}

\section{S2. Cryptographic Replay Audit Engine}
To avoid reliance on self-reported runtime metrics, the independent replay auditor (\nolinkurl{scripts/replay_and_score_p2c.py}) re-scores every Gazebo simulation run offline:

\subsection{S2.1 Replay Integrity Verification Steps}
\begin{enumerate}
    \item \textbf{Manifest Audit}: Verifies every serialized file against \texttt{checksums.sha256} to detect modifications relative to the recorded manifest (detecting file changes relative to the manifest, rather than asserting intrinsic authenticity of the original underlying sensor data).
    \item \textbf{Trajectory Distance}: Computes traversed distance by summing Euclidean distances between consecutive odometry sample positions $\langle x_k, y_k \rangle_{k=1}^K$:
    \begin{equation}
    d = \sum_{k=2}^K \sqrt{(x_k - x_{k-1})^2 + (y_k - y_{k-1})^2}
    \end{equation}
    \item \textbf{Halt Stability Verification}: Verifies that upon goal completion, the robot remains stationary over a minimum nominal window of $2.0\,\text{s}$:
    \begin{equation}
    \max_{t \in [t_{\text{goal}}, t_{\text{goal}} + \Delta t_{\text{stab}}]} |v_{\text{lin}}(t)| \le 0.05\,\text{m/s}, \quad \max_{t \in [t_{\text{goal}}, t_{\text{goal}} + \Delta t_{\text{stab}}]} |v_{\text{ang}}(t)| \le 0.08\,\text{rad/s}
    \end{equation}
    \item \textbf{Chronological Causality}:
    \begin{align}
    t_{\text{rec}} &\ge \max(t_{\text{action\_end}}, t_{\text{obs\_eval}}) - 0.05\,\text{s} \\
    t_{\text{inv}} &\ge \max(t_{\text{failure}}, t_{\text{free\_eval}}) - 0.05\,\text{s}
    \end{align}
\end{enumerate}

\section{S3. Reproducible PDF Build Pipeline}
The reference build pipeline (\texttt{paper/scripts/build\_paper\_pdf.py}):
\begin{enumerate}
    \item \textbf{Dynamic Table Injection}: Substitutes table markers with freshly generated LaTeX and Markdown tables.
    \item \textbf{MathJax Vector SVG Pre-Rendering}: Converts formulas into inline SVG vector paths via Node.js \texttt{mathjax-full@3.2.2}.
    \item \textbf{Table Row Validation}: Validates row counts of generated data tables prior to injection, distinguishing generated file names from paper table numbers: file \texttt{table1\_condition\_summary} (rendered as Table~2 in the main paper) requires exactly 10 condition rows; file \texttt{table2\_pairwise\_contrasts} (rendered as Table~3 in the main paper) requires exactly 10 contrast rows; and file \texttt{table3\_h1\_feasibility} (rendered as Table~4 in the main paper) requires exactly 4 trial rows (while Table~1 in the main paper presents the qualitative mechanism taxonomy).
    \item \textbf{Zero Raw TeX Leak Assertion}: Renders compiled PDF pages to PNG at $150\,\text{dpi}$ and extracts text (\texttt{pdftotext}) to ensure zero unrendered TeX control sequences.
    \item \textbf{Structured Build Record}: Emits \texttt{paper/build\_report.json} documenting tool versions, git status, and cryptographic digests.
\end{enumerate}

\section{S4. Detailed Parameter \& Source Code Traceability}
For exact code symbol and line-number mappings for all mathematical formulations and protocol parameters, refer to \texttt{paper/method-code-map.csv} (33 parameter mappings). For the complete literature primary verification ledger, refer to \texttt{paper/reference-verification.csv} (15 primary-source verified references).

\section{S5. D2 Trajectory Time Ledger: Observational Decomposition \& Threshold Sensitivity}

\subsection{S5.1 Observational Trajectory Time Ledger \& Algebraic Closure}
To analyze timing differences between FailMem ($F$) and persistent suppression ($M1$) in Scenario D2 without relying on unverified stationary assumptions, we perform an observational trajectory decomposition across Full Episode, History Phase, and Decision Phase.

\noindent\textbf{Notice on Algebraic Residuals}: \emph{The zero residual ($\text{residual} = 0.000\,\text{s}$) across all phases arises strictly by algebraic construction (absorbing unobserved active intervals and outside-action time as remainder components), and does not serve as independent validation of velocity classification accuracy or physical mechanism validity.}

\begin{table*}[t]
\centering
\footnotesize
\caption{Decision Phase Ledger ($t \ge T_{\text{history}}$).}
\label{tab:ledger_dec}
\setlength{\tabcolsep}{2.8pt}
\begin{tabular*}{\textwidth}{@{\extracolsep{\fill}}llccccccccc@{}}
\toprule
Condition & $n$ & \shortstack{Dist\\(m)} & \shortstack{Duration\\$T_{\text{dec}}$ (s)} & \shortstack{Moving\\$T_{\text{mov}}$ (s)} & \shortstack{Turning\\$T_{\text{turn}}$ (s)} & \shortstack{Stopped\\$T_{\text{stop}}$ (s)} & \shortstack{[Settling\\Window] (s)} & \shortstack{Unobserved\\$T_{\text{unobs}}$ (s)} & \shortstack{Outside\\Actions (s)} \\
\midrule
\textbf{D2\_F}  & 3 & 5.89 $\pm$ 0.02 & 51.10 $\pm$ 0.95 & 28.70 $\pm$ 1.18 & 4.20 $\pm$ 0.82 & 15.83 $\pm$ 2.42 & [6.97] & 1.33 $\pm$ 0.06 & 1.03 $\pm$ 0.15 \\
\textbf{D2\_M1} & 3 & 7.89 $\pm$ 0.02 & 51.30 $\pm$ 2.00 & 34.30 $\pm$ 1.76 & 3.47 $\pm$ 0.49 & 12.13 $\pm$ 1.41 & [4.60] & 0.70 $\pm$ 0.00 & 0.70 $\pm$ 0.00 \\
\textbf{D2\_O}  & 3 & 5.94 $\pm$ 0.07 & 52.03 $\pm$ 1.75 & 29.33 $\pm$ 2.15 & 3.83 $\pm$ 0.49 & 16.47 $\pm$ 0.06 & [6.93] & 1.30 $\pm$ 0.10 & 1.10 $\pm$ 0.20 \\
\textbf{D2\_R}  & 3 & 5.92 $\pm$ 0.04 & 49.57 $\pm$ 1.59 & 28.27 $\pm$ 0.45 & 4.17 $\pm$ 0.84 & 14.70 $\pm$ 1.55 & [5.37] & 1.27 $\pm$ 0.06 & 1.17 $\pm$ 0.12 \\
\bottomrule
\end{tabular*}
\end{table*}

\begin{table*}[t]
\centering
\footnotesize
\caption{History Phase Ledger ($t < T_{\text{history}}$).}
\label{tab:ledger_hist}
\setlength{\tabcolsep}{2.8pt}
\begin{tabular*}{\textwidth}{@{\extracolsep{\fill}}llccccccccc@{}}
\toprule
Condition & $n$ & \shortstack{Dist\\(m)} & \shortstack{Duration\\$T_{\text{hist}}$ (s)} & \shortstack{Moving\\$T_{\text{mov}}$ (s)} & \shortstack{Turning\\$T_{\text{turn}}$ (s)} & \shortstack{Stopped\\$T_{\text{stop}}$ (s)} & \shortstack{[Settling\\Window] (s)} & \shortstack{Unobserved\\$T_{\text{unobs}}$ (s)} & \shortstack{Outside\\Actions (s)} \\
\midrule
\textbf{D2\_F}  & 3 & 10.36 $\pm$ 0.08 & 100.00 $\pm$ 7.72 & 52.37 $\pm$ 3.66 & 9.93 $\pm$ 0.75 & 30.13 $\pm$ 4.31 & [11.00] & 3.17 $\pm$ 0.15 & 4.40 $\pm$ 0.30 \\
\textbf{D2\_M1} & 3 & 10.35 $\pm$ 0.11 & 97.53 $\pm$ 2.12  & 49.23 $\pm$ 1.58 & 10.10 $\pm$ 0.76 & 30.63 $\pm$ 0.91 & [11.77] & 3.20 $\pm$ 0.10 & 4.37 $\pm$ 0.31 \\
\textbf{D2\_O}  & 3 & 10.46 $\pm$ 0.22 & 97.73 $\pm$ 1.21  & 49.80 $\pm$ 1.14 & 9.83 $\pm$ 0.50  & 30.70 $\pm$ 0.78 & [11.87] & 3.23 $\pm$ 0.06 & 4.17 $\pm$ 0.42 \\
\textbf{D2\_R}  & 3 & 10.43 $\pm$ 0.05 & 98.60 $\pm$ 2.38  & 50.90 $\pm$ 1.32 & 9.40 $\pm$ 0.46  & 30.53 $\pm$ 1.80 & [11.77] & 3.07 $\pm$ 0.06 & 4.70 $\pm$ 0.27 \\
\bottomrule
\end{tabular*}
\end{table*}

\begin{table*}[t]
\centering
\footnotesize
\caption{Full Episode Ledger.}
\label{tab:ledger_full}
\setlength{\tabcolsep}{2.8pt}
\begin{tabular*}{\textwidth}{@{\extracolsep{\fill}}llccccccccc@{}}
\toprule
Condition & $n$ & \shortstack{Dist\\(m)} & \shortstack{Duration\\$T_{\text{total}}$ (s)} & \shortstack{Moving\\$T_{\text{mov}}$ (s)} & \shortstack{Turning\\$T_{\text{turn}}$ (s)} & \shortstack{Stopped\\$T_{\text{stop}}$ (s)} & \shortstack{[Settling\\Window] (s)} & \shortstack{Unobserved\\$T_{\text{unobs}}$ (s)} & \shortstack{Outside\\Actions (s)} \\
\midrule
\textbf{D2\_F}  & 3 & 16.25 $\pm$ 0.10 & 151.10 $\pm$ 7.76 & 81.07 $\pm$ 4.48 & 14.13 $\pm$ 1.29 & 45.97 $\pm$ 2.57 & [17.97] & 4.50 $\pm$ 0.17 & 5.43 $\pm$ 0.45 \\
\textbf{D2\_M1} & 3 & 18.25 $\pm$ 0.09 & 148.83 $\pm$ 2.40 & 83.53 $\pm$ 3.01 & 13.57 $\pm$ 0.51 & 42.77 $\pm$ 1.61 & [16.37] & 3.90 $\pm$ 0.10 & 5.07 $\pm$ 0.31 \\
\textbf{D2\_O}  & 3 & 16.40 $\pm$ 0.24 & 149.77 $\pm$ 0.55 & 79.13 $\pm$ 1.03 & 13.67 $\pm$ 0.15 & 47.17 $\pm$ 0.76 & [18.80] & 4.53 $\pm$ 0.12 & 5.27 $\pm$ 0.50 \\
\textbf{D2\_R}  & 3 & 16.34 $\pm$ 0.03 & 148.17 $\pm$ 3.16 & 79.17 $\pm$ 0.90 & 13.57 $\pm$ 1.21 & 45.23 $\pm$ 1.16 & [17.13] & 4.33 $\pm$ 0.06 & 5.87 $\pm$ 0.38 \\
\bottomrule
\end{tabular*}
\end{table*}

\noindent\emph{Note}: Settling window duration (\texttt{sim\_duration\_covered}) is reported strictly as an independent event-level observation metric and is not algebraically subtracted from stopped time.

\subsection{S5.2 Discretization Threshold Sensitivity Analysis}
To evaluate sensitivity to kinematic classification parameters, Decision Phase moving and stopped durations were evaluated across linear velocity thresholds $v_{\text{lin}} \in \{0.03, 0.05, 0.08\}\,\text{m/s}$, angular velocity thresholds $\omega_{\text{ang}} \in \{0.05, 0.08, 0.12\}\,\text{rad/s}$, and gap thresholds $\tau_{\text{gap}} \in \{0.25, 0.50, 1.00\}\,\text{s}$.

\begin{table}[h]
\centering
\footnotesize
\caption{Sensitivity of Moving and Stopped Durations to Discretization Thresholds.}
\label{tab:sensitivity}
\setlength{\tabcolsep}{4pt}
\begin{tabular}{ccccccccc}
\toprule
$v_{\text{lin}}$ (m/s) & $\omega_{\text{ang}}$ (rad/s) & $\tau_{\text{gap}}$ (s) & $F$ Moving (s) & $M1$ Moving (s) & Diff (s) & $F$ Stopped (s) & $M1$ Stopped (s) & Diff (s) \\
\midrule
0.03 & 0.05 & 0.25 & 31.30 & 36.40 & -5.10 & 13.80 & 10.70 & +3.10 \\
0.03 & 0.08 & 0.50 & 31.30 & 36.40 & -5.10 & 13.80 & 10.70 & +3.10 \\
0.03 & 0.12 & 1.00 & 31.30 & 36.40 & -5.10 & 13.80 & 10.70 & +3.10 \\
0.05 & 0.05 & 0.25 & 28.70 & 34.30 & -5.60 & 15.77 & 12.07 & +3.70 \\
0.05 & 0.08 & 0.50 & 28.70 & 34.30 & -5.60 & 15.83 & 12.13 & +3.70 \\
0.05 & 0.12 & 1.00 & 28.70 & 34.30 & -5.60 & 15.90 & 12.13 & +3.77 \\
0.08 & 0.05 & 0.25 & 25.80 & 32.60 & -6.80 & 18.33 & 13.13 & +5.20 \\
0.08 & 0.08 & 0.50 & 25.80 & 32.60 & -6.80 & 18.47 & 13.27 & +5.20 \\
0.08 & 0.12 & 1.00 & 25.80 & 32.60 & -6.80 & 18.53 & 13.30 & +5.23 \\
\bottomrule
\end{tabular}
\end{table}

\noindent\textbf{Sensitivity Summary}: Across thresholds, sample mean moving differences range from $-5.10\,\text{s}$ to $-6.80\,\text{s}$, and stopped differences range from $+3.10\,\text{s}$ to $+5.23\,\text{s}$. Because individual components depend on discretization thresholds, we report them as observational latencies without asserting exact physical compensation.

\end{document}
